% !TeX root = ../main.tex

% ==================================================
% D.3. 자료형·Nullable·결측 표현
% ==================================================

\subsection{자료형·Nullable·결측 표현}

\subsubsection{이전 학습과의 연결}

\zsh{D.2}에서는 Pandas의 세 가지 핵심 객체인 Series, DataFrame, Index를 분석했다.

특히 다음과 같은 구조를 확인했다.

\begin{dhkcode}
import pandas as pd

df = pd.DataFrame({
    "text": ["hello", "world"],
    "label": [1, 0],
})

print(df.dtypes)
\end{dhkcode}

예상 출력:

\begin{dhkzsh}
text     object
label     int64
dtype: object
\end{dhkzsh}

여기서 \zsh{df.dtypes}는 각 열의 자료형을 나타내는 Series를 반환한다.

그런데 데이터가 다음처럼 바뀌면 어떻게 될까?

\begin{dhkcode}
df = pd.DataFrame({
    "text": ["hello", None],
    "label": [1, None],
})
\end{dhkcode}

문자열 열과 정수 열 모두에 \zsh{None}을 넣었지만, 두 열의 자료형은 같은 방식으로 결정되지 않을 수 있다.

이유는 Pandas가 단순히 Python 객체를 나열하는 시스템이 아니라, 열의 값들을 특정 자료형과 저장 표현에 따라 관리하는 시스템이기 때문이다.

이번 학습에서는 이 구조를 정확히 분석할 것이다.

\subsubsection{이번 학습의 핵심 질문}

이번 \zsh{D.3}에서는 다음 다섯 가지 질문에 답해야 한다.

\begin{dhksmry}
1. Python 객체의 type과 Pandas의 dtype은 어떻게 다른가?

2. int64와 Int64, bool과 boolean, object와 string은 왜 구분되는가?

3. None, np.nan, pd.NA, pd.NaT는 각각 무엇을 표현하는가?

4. 결측값이 포함된 비교·논리·산술 연산은 어떤 결과를 반환하는가?

5. ML/Transformer 데이터셋에서 결측을 잘못 처리하면 왜 정답 레이블과
   평과 결과가 오염되는가?
\end{dhksmry}

이 질문들은 모두 하나의 원리로 연결된다.

데이터 값의 의미, 그 값을 저장하는 자료형, 그리고 결측 상태를 표현하는 방법은 서로 구분해야 한다는 원리다.

\subsubsection{\zsh{type}과 \zsh{dtype}은 서로 다른 계층의 정보다}

Python에서 \zsh{type(obj)}는 객체의 타입을 반환한다.

\begin{dhkcode}
x = 10

print(type(x))
\end{dhkcode}

예상 출력:

\begin{dhkzsh}
<class 'int'>
\end{dhkzsh}

여기서 \zsh{x}는 Python의 \zsh{int} 객체를 참조한다.

반면 Pandas의 \zsh{Series}는 여러 값을 관리하는 하나의 객체다.

\begin{dhkcode}
import pandas as pd

s = pd.Series([10, 20, 30], dtype="int64")

print(type(s))
print(s.dtype)
print(type(s.iloc[0]))
\end{dhkcode}

일반적인 NumPy 기반 \zsh{int64} Series에서 예상되는 출력은 다음과 같다.

\begin{dhkzsh}
<class 'pandas.core.series.Series'>
int64
<class 'numpy.int64'>
\end{dhkzsh}

세 결과가 각각 다른 계층을 나타낸다는 점에 주목해야 한다.

\begin{dhktable}
|표현식|	확인하는 대상|	의미|
|:---|:---|:---|
|type(s)	|Series 객체|	어떤 Python 클래스의 인스턴스인가?|
|s.dtype	|Series의 값 저장 자료형|	값들이 어떤 dtype으로 표현되는가?|
|type(s.iloc[0])|	선택된 스칼라 객체|	조회 결과가 어떤 Python/NumPy 스칼라인가?|
\end{dhktable}

\zsh{dtype}은 Series 객체의 Python 클래스가 아니다.

Series가 관리하는 데이터의 자료형을 기술하는 정보이다.

또한 \zsh{dtype="int64"}라고 해서 각 원소가 Python의 내장 \zsh{int} 객체로 개별 저장된다는 뜻도 아니다.

NumPy 기반 정수 배열에서는 정수들이 고정 폭의 데이터 표현으로 저장되고, 원소를 조회할 때 NumPy 스칼라 객체를 얻을 수 있다.

객체 타입과 객체 내부 상태의 구분을 Pandas에 적용할 사례라고 볼 수 있다.

\begin{dhknote}
\begin{dhksmry}[centered]
Pandas 자료형의 두 가지 주요 계열
\end{dhksmry}
\end{dhknote}

Pandas에서는 크게 다음 두 계열을 구분하는 것이 중요하다.

\begin{dhktable}
|계열|	예시|	특징|
|:---|:---|:---|
|NumPy 기반 dtype|	int64, float64, bool, datetime64[ns]|	NumPy 자료형 체계를 활용|
|Pandas ExtensionDtype|	Int64, Float64, boolean, string, category|	Pandas가 확장한 자료형 및 연산 규칙 제공|
\end{dhktable}

여기서 주의할 점은 ExtensionDtype이라고 해서 항상 별도의 결측 마스크를 사용하는 것은 아니라는 사실이다.

실제 저장 구조는 dtype과 배열 구현에 따라 달라진다.

예를 들어 Nullable 정수 배열은 일반적으로 데이터 배열과 결측 여부를 나타내는 마스크를 함께 관리하지만, 범주형과 문자열 확장 배열의 내부 구현까지 같은 구조라고 가정할 수는 없다.

이제 가장 중요한 Nullable 정수형부터 살펴본다.

\subsubsection{정수 열에 결측값이 들어가면 왜 \zsh{float64}가 될까?}

먼저 가장 기본적인 예제를 살펴보자.

\begin{dhknote}
\begin{dhksmry}[centered]
독립 실행 코드 1 - NumPy 정수형과 Nullable 정수형
\end{dhksmry}
\end{dhknote}

\begin{dhkcode}
import pandas as pd

a = pd.Series([1, None])
b = pd.Series([1, None], dtype="Int64")

print(a.dtype)
print(a.tolist())

print(b.dtype)
print(b.tolist())

print(type(a.iloc[0]).__name__)
print(type(b.iloc[0]).__name__)
\end{dhkcode}

예상 출력:

\begin{dhkzsh}
float64
[1.0, nan]
Int64
[1, <NA>]
float64
int64
\end{dhkzsh}

같은 값을 전달했는데 결과가 달라졌다.

왜 그럴까?

\begin{dhknote}
\begin{dhksmry}[centered]
1. a의 생성 과정
\end{dhksmry}
\end{dhknote}

\begin{dhkcode}
a = pd.Series([1, None])
\end{dhkcode}

생성자에 \zsh{dtype}을 지정하지 않았으므로 Pandas는 입력값으로부터 자료형을 추론한다.

NumPy의 일반적인 \zsh{int64} 배열에는 \zsh{None}을 정수 결측값으로 표현하는 규칙이 없다.

따라서 이 입력에 대해 Pandas는 \zsh{float64}로 변환하고, 결측 위치에 \zsh{NaN}을 사용한다.

결과적으로 정수 1도 실수 1.0으로 표현된다.

\begin{dhkbranch}[in=8]
node main :: 입력 // [1, None]
arrow main
node main :: dtype 추론
arrow main
node main :: float64 Series // [1.0, NaN]
\end{dhkbranch}

여기서 중요한 사실은 정수 데이터에 결측값이 추가되면서 값의 표현 자료형이 변경되었다는 점이다.

\begin{dhknote}
\begin{dhksmry}[centered]
2. b의 생성 과정
\end{dhksmry}
\end{dhknote}

\begin{dhkcode}
b = pd.Series([1, None], dtype="Int64")
\end{dhkcode}

이번에는 Pandas의 Nullable 정수형인 \zsh{Int64}를 명시했다.

대문자 \zsh{I}에 주목해야 한다.

\begin{dhksmry}
    - int64 : 일반적인 NumPy 기반 정수 dtype

    - Int64 : Pandas의 Nullable 정수 ExtensionDtype
\end{dhksmry}

\zsh{Int64}는 정수값과 결측 상태를 함께 표현할 수 있다.

따라서 정수 1을 실수로 변환하지 않고도 두 번째 위치가 결측이라는 사실을 표현한다.

\begin{dhknote}
\begin{dhksmry}[centered]
Int64의 내부 구조
\end{dhksmry}
\end{dhknote}

Nullable 정수형의 논리적 구조는 다음과 같이 이해할 수 있다.

\begin{dhkgraph}[uniform]
node n1 (1,1) :: Series 객체 b // dtype: Int64
node n2 (1,2) :: 데이터 영역 // 1 ?
node n3 (1,3) :: 결측 마스크 // F T
group g1(n2,n3) :: IntegerArray

edge n1 -> g1

explain :: 마스크가 True인 위치는 실제 데이터 영역의 값과 관계없이 결측으로 해석

caption :: Nullable 정수 배열의 개념적 구조. // ? 는 결측 위치의 내부 저장값을 특정하지 않는다는 뜻이다.
\end{dhkgraph}

실제 객체 타입도 확인할 수 있다.

\begin{dhkcode}
import pandas as pd

s = pd.Series([10, None, 30], dtype="Int64")

print(type(s.array).__name__)
print(s.array.dtype)
print(s.isna().tolist())
\end{dhkcode}

예상 출력:

\begin{dhkzsh}
IntegerArray
Int64
[False, True, False]
\end{dhkzsh}

\zsh{Series.array}는 Series가 관리하는 데이터의 배열 인터페이스를 제공한다.

이 경우 반환 객체는 \zsh{IntegerArray}다.

\zsh{IntegerArray}는 데이터와 결측 마스크를 관리할 수 있지만, 내부의 비공개 속성인 \zsh{_data}, \zsh{_mask} 등에 직접 의존하는 코드는 피해야 한다.

그 구조는 공개 API의 안정적인 계약이 아니기 때문이다.

이전 학습과 연결하면 다음과 같이 구분할 수 있다.

\begin{dhksmry}
    - s는 Series 인스턴스를 참조하는 이름이다.

    - s.array는 Series의 배열 객체를 조회하는 속성 접근이다.

    - s.dtype은 해당 배열의 자료형을 기술하는 객체다.

    - s.isna()는 결측 여부를 계산하여 새로운 불리언 Series를
      반환하는 메서드 호출이다.
\end{dhksmry}

결측 상태가 있다는 사실과 값 자체가 특정 숫자라는 사실은 독립적으로 관리될 수 있다.

\subsubsection{\zsh{None}, \zsh{np.nan}, \zsh{pd.NA}, \zsh{pd.NaT}는 무엇이 다를까?}

Pandas에서 결측을 표현하는 대표적인 객체는 네 가지다.

\begin{dhktable}
|객체|	주된 의미|	일반적인 사용 맥락|
|None|	Python의 값 부재를 나타내는 싱글턴 객체|	일반 Python 객체, object 열|
|np.nan|	IEEE 부동소수점 NaN 값|	NumPy 실수 배열, 일반 실수 Series|
|pd.NA|	Pandas의 결측 스칼라|	Nullable 정수·불리언·문자열 등|
|pd.NaT|	Not a Time|	날짜·시간·기간 데이터의 결측|
\end{dhktable}

이들은 모두 결측을 표현할 수 있지만, 동일한 객체도 아니고 비교 연산의 의미도 같지 않다.

\begin{dhknote}
\begin{dhksmry}[centered]
독립 실행 코드 2 - 객체 타입과 비교
\end{dhksmry}
\end{dhknote}

\begin{dhkcode}
import numpy as np
import pandas as pd

values = {
    "None": None,
    "NaN": np.nan,
    "NA": pd.NA,
    "NaT": pd.NaT,
}

for name, value in values.items():
    print(
        name,
        type(value).__name__,
        pd.isna(value),
        value == value,
    )
\end{dhkcode}

예상 출력:

\begin{dhkzsh}
None NoneType True True
NaN float True False
NA NAType True <NA>
NaT NaTType True False
\end{dhkzsh}

여기에는 매우 중요한 차이가 있다.

\begin{dhknote}
\begin{dhksmry}[centered]
1. None == None은 True
\end{dhksmry}
\end{dhknote}

\zsh{None}은 Python에서 유일한 싱글턴 객체다.

\begin{dhkcode}
print(None is None)
print(None == None)
\end{dhkcode}

예상 출력:

\begin{dhkzsh}
True
True
\end{dhkzsh}

따라서 \zsh{None}의 동일성과 동등성 비교는 모두 참이다.

하지만 이것은 \zsh{None}이 결측이 아니라는 뜻이 아니다.

Python 객체의 비교 규칙과 데이터 분석에서 결측으로 해석하는 규칙은 서로 다른 계층에 있다.

\begin{dhknote}
\begin{dhksmry}[centered]
2. np.nan == np.nan은 False
\end{dhksmry}
\end{dhknote}

\zsh{NaN}은 일반적인 숫자와 다른 IEEE 부동소수점 비교 규칙을 따른다.

\begin{dhkcode}
import numpy as np

print(np.nan == np.nan)
\end{dhkcode}

예상 출력:

\begin{dhkzsh}
False
\end{dhkzsh}

따라서 \zsh{x == np.nan}을 결측 검사 방법으로 사용해서는 안 된다.

\begin{dhknote}
\begin{dhksmry}[centered]
3. pd.NA == pd.NA는 <NA>
\end{dhksmry}
\end{dhknote}

\begin{dhkcode}
import pandas as pd

print(pd.NA == pd.NA)
\end{dhkcode}

예상 출력:

\begin{dhkzsh}
<NA>
\end{dhkzsh}

이것은 결측값끼리 비교한 결과를 참이나 거짓으로 확정할 수 없다는 의미를 표현한다.

예를 들어 어떤 학생의 점수가 누락되어 있다면, 그 점수가 다른 누락된 점수와 같다고 단정할 수 없을 것이다.

Pandas의 Nullable 연산은 이러한 불확정성을 표현할 수 있다.

다만 모든 연산이 예외 없이 \zsh{<NA>}를 전파하는 것은 아니다.

특히 불리언 연산에서는 다른 피연산자의 값에 따라 결과를 확정할 수 있다.

이 부분은 \zsh{D.3.7}에서 살펴볼 것이다.

\begin{dhknote}
\begin{dhksmry}[centered]
결측 검사는 pd.isna()를 사용한다
\end{dhksmry}
\end{dhknote}

\begin{dhkcode}
import numpy as np
import pandas as pd

print(pd.isna(None))
print(pd.isna(np.nan))
print(pd.isna(pd.NA))
print(pd.isna(pd.NaT))
\end{dhkcode}

예상 출력:

\begin{dhkzsh}
True
True
True
True
\end{dhkzsh}

\zsh{pd.isna()}는 Pandas가 인식하는 결측 상태를 검사하는 함수다.

스칼라를 전달하면 불리언 스칼라를, Series를 전달하면 각 원소의 결측 여부를 나타내는 Series를 반환할 수 있다.

따라서 객체의 동등성 비교와 결측 여부 검사는 서로 다른 연산이라는 점을 기억해야 한다.

\subsubsection{Nullable 자료형은 왜 별도로 필요할까?}

\zsh{Int64}뿐 아니라 Pandas에는 여러 Nullable 자료형이 있다.

\begin{dhktable}
|일반적인 dtype|	Nullable 대응 dtype|	주요 용도|
|:---|:---|:---|
|int64|	Int64|	누락 가능한 정수|
|float64|	Float64|	누락 가능한 실수|
|bool|	boolean|	참·거짓·결측|
|object|	string|	누락 가능한 문자열 데이터|
\end{dhktable}

여기서 \zsh{object}는 반드시 문자열만 저장한다는 뜻이 아니다.

Python 객체를 참조하는 범용 자료형이기 때문에 문자열, 숫자, 리스트 등 서로 다른 종류의 객체가 들어갈 수도 있다.

\begin{dhknote}
\begin{dhksmry}[centered]
독립 실행 코드 3 - 자료형 추론과 명시적 변환
\end{dhksmry}
\end{dhknote}

\begin{dhkcode}
import pandas as pd

df = pd.DataFrame({
    "count": [1, None, 3],
    "active": [True, None, False],
    "text": ["ko", None, "en"],
})

print(df.dtypes)

converted = df.convert_dtypes()

print()
print(converted.dtypes)
\end{dhkcode}

Pandas 2.2.3 기준 실행 결과:

\begin{dhkzsh}
count     float64
active     object
text       object
dtype: object

count              Int64
active           boolean
text      string[python]
dtype: object
\end{dhkzsh}

\zsh{convert_dtypes()}는 기존 데이터에 적합한 Nullable 자료형을 추론해 새로운 DataFrame을 반환하는 메서드다.

여기서 주의할 점이 있다.

\begin{dhkcode}
converted = df.convert_dtypes()
\end{dhkcode}

이 코드는 \zsh{df}라는 이름을 새로운 객체에 재바인딩하는 것이 아니다.

\zsh{convert_dtypes()}의 반환값을 \zsh{converted}라는 새로운 이름에 바인딩한 것이다.

원본 \zsh{df}는 그대로 유지된다.

또한 \zsh{convert_dtypes()}는 데이터를 분석해 dtype을 추론하는 편의 기능이지, 모든 열의 업무상 의미를 판별하는 검증기가 아니다.

예를 들어 문자열 \zsh{"00123"}이 우편번호라면 숫자로 바꾸는 것이 오히려 정보 손실을 일으킬 수 있을 것이다.

자료형 추론과 데이터 의미 검증은 구분해야 한다.

\subsubsection{Nullable Boolean과 3값 논리}

일반적인 Python의 불리언은 \zsh{True}와 \zsh{False} 두 가지 값으로 표현된다.

하지만 데이터 분석에서는 조건의 참·거짓을 판정할 정보가 없는 경우도 존재한다.

Pandas의 \zsh{boolean} dtype은 다음 세 상태를 표현할 수 있다.

\begin{dhksmry}
    - True

    - False

    - <NA>
\end{dhksmry}

이것은 \zsh{3값 논리(three-valued logic)}와 연결된다.

\begin{dhknote}
\begin{dhksmry}[centered]
pd.NA를 Python bool로 변환하면?
\end{dhksmry}
\end{dhknote}

\begin{dhkcode}
import pandas as pd

print(bool(pd.NA))
\end{dhkcode}

이 코드는 다음 오류를 발생시킨다.

\begin{dhkzsh}
TypeError: boolean value of NA is ambiguous
\end{dhkzsh}

\zsh{pd.NA}가 참인지 거짓인지 결정할 수 없기 때문이다.

그런데 불리언 연산에서는 일부 결과를 확정할 수 있다.

\begin{dhknote}
\begin{dhksmry}[centered]
독립 실행 코드 4 - 3값 논리
\end{dhksmry}
\end{dhknote}

\begin{dhkcode}
import pandas as pd

print(pd.NA & False)
print(pd.NA & True)

print(pd.NA | True)
print(pd.NA | False)
\end{dhkcode}

예상 출력:

\begin{dhkzsh}
False
<NA>
True
<NA>
\end{dhkzsh}

왜 이런 결과가 발생할까?

\begin{dhktable}
|연산|	결과	|이유|
|:---|:---|:---|
|<NA> & False|	False|	다른 값이 무엇이든 AND 결과는 거짓|
|<NA> & True|	<NA>|	결측값에 따라 결과가 달라짐|
|<NA> ㅣ True|	True|	다른 값이 무엇이든 OR 결과는 참|
|<NA> ㅣ False|	<NA>|	결측값에 따라 결과가 달라짐|
\end{dhktable}

이것은 \zsh{Kleene} 논리라고 부른다.

여기서 \zsh{&}, \zsh{|}는 Python의 비트 연산자 문법이지만 Pandas 객체에서는 원소별 논리 연산을 수행하도록 구현되어 있다.

따라서 \zsh{and}, \zsh{or}와 혼동하면 안 된다.

특히 \zsh{and}와 \zsh{or}는 피연산자의 진리값을 평가하는 Python 언어 연산자이고, Pandas Series의 원소별 논리 연산을 대신하지 않는다.

\begin{dhknote}
\begin{dhksmry}[centered]
결측이 포함된 비교 연산
\end{dhksmry}
\end{dhknote}

\begin{dhkcode}
import pandas as pd

s = pd.Series([1, None], dtype="Int64")

result = s == 1

print(result.dtype)
print(result.tolist())
print(result.fillna(False).tolist())
\end{dhkcode}

예상 출력:

\begin{dhkzsh}
boolean
[True, <NA>]
[True, False]
\end{dhkzsh}

\zsh{result}는 Nullable Boolean Series다.

여기서 \zsh{fillna(False)}는 단순한 출력 형식 변경이 아니다.

판정할 수 없는 상태를 거짓으로 처리한다는 정책을 적용한 것이다.

이 정책은 분석 목적에 따라 적절할 수도, 부적절할 수도 있다.

예를 들어 \zsh{"정답 레이블이 1인가?"}라는 조건에서 결측 레이블을 \zsh{False}로 바꾸면, 실제로는 정답이 알려지지 않은 샘플을 클래스 1이 아닌 샘플처럼 해석할 위험이 있다.

참고로 Nullable Boolean 마스크를 행 선택에 사용할 때 결측 위치는 일반적으로 선택되지 않는 것으로 처리된다.

그러나 행 선택에서 제외된다는 사실이 원래 조건의 값이 \zsh{False}였다는 뜻은 아니다.

이 차이는 \zsh{D.6}에서 다시 분석할 것이다.

\subsubsection{문자열·범주·날짜 자료형}

이제 정수와 불리언 이외의 자료형을 살펴보자.

\begin{dhknote}
\begin{dhksmry}[centered]
1. 문자열: object와 string
\end{dhksmry}
\end{dhknote}

\begin{dhkcode}
import pandas as pd

s = pd.Series(
    ["ko", "en", None],
    dtype="string[python]",
)

print(s.dtype)
print(s.tolist())
\end{dhkcode}

예상 출력:

\begin{dhkzsh}
string
['ko', 'en', <NA>]
\end{dhkzsh}

\zsh{string[python]}은 Pandas의 문자열 확장 자료형을 Python 기반 저장 방식으로 명시한 것이다.

\zsh{object} dtype은 문자열 외의 Python 객체도 저장할 수 있지만, \zsh{string} dtype은 문자열 데이터라는 의미를 더 명확하게 표현한다.

Pandas 버전에 따라 문자열 dtype의 기본 추론 및 저장 백엔드가 달라질 수 있으므로, 재현성이 중요한 예제에서는 \zsh{string[python]}처럼 명시하는 것이 좋다.

\begin{dhknote}
\begin{dhksmry}[centered]
2. 범주형: category
\end{dhksmry}
\end{dhknote}

범주형 자료형은 값들이 제한된 범주 집합에 속한다는 정보를 표현한다.

\begin{dhkcode}
import pandas as pd

s = pd.Series(
    ["train", "valid", None],
    dtype="category",
)

print(s.dtype)
print(s.cat.categories.tolist())
print(s.isna().tolist())
\end{dhkcode}

예상 출력:

\begin{dhkzsh}
category
['train', 'valid']
[False, False, True]
\end{dhkzsh}

범주형은 실제 범주값과 내부 코드의 대응 관계를 이용할 수 있다.

이 자료형의 중요한 특징은 단순히 문자열을 저장하는 것과 달리, 가능한 범주와 범주의 순서를 자료형의 일부로 표현할 수 있다는 점이다.

다만 \zsh{category}가 자동으로 업무상 올바른 범주를 정의해 주는 것은 아니다.

범주 집합과 순서가 중요하다면 \zsh{CategoricalDtype}으로 명시해야 한다.

\begin{dhknote}
\begin{dhksmry}[centered]
3. 날짜·시간: datetime64[ns]와 NaT
\end{dhksmry}
\end{dhknote}

\begin{dhkcode}
import pandas as pd

s = pd.to_datetime(
    pd.Series([
        "2026-11-01",
        "invalid",
        None,
    ]),
    errors="coerce",
)

print(s.dtype)
print(s.isna().tolist())
\end{dhkcode}

예상 출력:

\begin{dhkzsh}
datetime64[ns]
[False, True, True]
\end{dhkzsh}

\zsh{errors="coerce"}는 변환할 수 없는 값을 날짜·시간 결측값인 \zsh{NaT}로 바꾸도록 지시한다.

여기서 한 가지 중요한 문제가 발생한다.

원래부터 비어 있던 값과 날짜 형식이 잘못되어 변환에 실패한 값이 모두 \zsh{NaT}가 될 수 있다.

따라서 변환 결과만 보관하면 원래 누락된 데이터와 파싱 실패 데이터를 구분할 수 없게 될 수 있다.

이 문제는 ML 데이터셋 품질 검증에서 중요하므로 \zsh{D.3.11}에서 다시 다룰 것이다.

\subsubsection{자료형 변환과 오류 발생 원리}

Pandas에서 자료형 변환에 사용하는 대표적인 메서드는 \zsh{astype()}이다.

\begin{dhkcode}
s.astype("Int64")
\end{dhkcode}

이 메서드는 변환된 자료형의 Series를 반환한다.

\begin{dhknote}
\begin{dhksmry}[centered]
독립 실행 코드 5 - 변환 실패와 결측 처리
\end{dhksmry}
\end{dhknote}

\begin{dhkcode}
import pandas as pd

raw = pd.Series(
    ["12", "bad", None],
    dtype="string[python]",
)

try:
    converted = raw.astype("Int64")
except ValueError as e:
    print(type(e).__name__)

parsed = pd.to_numeric(
    raw,
    errors="coerce",
).astype("Int64")

print(parsed.tolist())

invalid = raw.notna() & parsed.isna()

print(invalid.tolist())
\end{dhkcode}

예상 출력:

\begin{dhkzsh}
ValueError
[12, <NA>, <NA>]
[False, True, False]
\end{dhkzsh}

실행 과정을 분석해 보자.

\begin{dhknote}
\begin{dhksmry}[centered]
첫 번째 변환
\end{dhksmry}
\end{dhknote}

\begin{dhkcode}
raw.astype("Int64")
\end{dhkcode}

\zsh{"12"}는 정수로 변환할 수 있지만 \zsh{"bad"}는 정수로 해석할 수 없다.

따라서 변환이 실패하고 예외가 발생한다.

\begin{dhknote}
\begin{dhksmry}[centered]
두 번째 변환
\end{dhksmry}
\end{dhknote}

\begin{dhkcode}
pd.to_numeric(raw, errors="coerce")
\end{dhkcode}

이번에는 변환할 수 없는 값을 결측으로 처리하도록 명시했다.

따라서 \zsh{"bad"}와 \zsh{None}이 모두 결측이 된다.

\begin{dhknote}
\begin{dhksmry}[centered]
세 번째 단계
\end{dhksmry}
\end{dhknote}

\begin{dhkcode}
invalid = raw.notna() & parsed.isna()
\end{dhkcode}

원래 값이 존재했는데 변환 결과가 결측인 위치를 찾는다.

그 결과 두 번째 원소만 \zsh{True}가 된다.

이것은 단순한 dtype 변환보다 한 단계 더 중요한 데이터 처리 원칙을 보여준다.

결측값으로 변환하기 전에, 결측이 발생한 원인을 보존해야 할 수 있다.

실전 데이터에서는 다음 상태를 구분해야 하는 경우가 많다.

\begin{dhktable}
|상태	|예시|	의미|
|:---|:---|:---|
|원래 누락|	None	|입력값이 존재하지 않음|
|파싱 실패|	"bad"	|입력값은 있지만 기대한 형식이 아님|
|적용 불가|	해당 없음	|해당 항목 자체가 의미 없음|
|아직 미확정|	레이블 검수 중|	추후 값이 결정될 수 있음|
\end{dhktable}

이들은 모두 결측값으로 표현될 수 있지만, 같은 업무상 의미를 갖는 것은 아니다.

\subsubsection{결측값과 집계 연산}

결측값의 표현 방식은 산술 연산과 집계 결과에도 영향을 준다.

\begin{dhknote}
\begin{dhksmry}[centered]
독립 실행 코드 6 - 비교와 합계
\end{dhksmry}
\end{dhknote}

\begin{dhkcode}
import numpy as np
import pandas as pd

a = pd.Series([1.0, np.nan])

b = pd.Series(
    [1, pd.NA],
    dtype="Int64",
)

print((a == a).tolist())
print((b == b).tolist())

print(a.sum())
print(b.sum())

print(a.sum(skipna=False))
print(b.sum(skipna=False))
\end{dhkcode}

예상 출력:

\begin{dhkzsh}
[True, False]
[True, <NA>]
1.0
1
nan
<NA>
\end{dhkzsh}

여기서 두 가지 규칙을 구분해야 한다.

첫째, 원소별 비교 연산은 dtype의 결측 표현과 비교 의미론을 따른다는 점이다.

일반적인 \zsh{float64}에서 \zsh{NaN}은 자기 자신과 비교해도 \zsh{False}를 반환한다.

반면 Nullable 정수형에서는 결측값의 동등성을 판정할 수 없으므로 \zsh{<NA>}가 반환된다.

둘째, 집계 메서드는 결측을 건너뛸지 결정하는 별도의 정책을 갖는다는 점이다.

\zsh{sum()}은 기본적으로 결측을 건너뛰는 \zsh{skipna=True} 동작을 사용한다.

따라서 두 Series 모두 관측된 값 1만 합산한다.

반면 \zsh{skipna=False}를 전달하면 결측을 무시하지 않으므로 결과도 결측이 된다.

\begin{dhknote}
\begin{dhksmry}[centered]
모든 값이 결측이면?
\end{dhksmry}
\end{dhknote}

여기에는 추가적인 경계 조건이 있다.

\begin{dhkcode}
import pandas as pd

s = pd.Series(
    [pd.NA, pd.NA],
    dtype="Int64",
)

print(s.sum())
print(s.sum(min_count=1))
print(s.count())
\end{dhkcode}

예상 출력:

\begin{dhkzsh}
0
<NA>
0
\end{dhkzsh}

왜 첫 번째 결과가 0일까?

\zsh{sum()}은 기본적으로 \zsh{min_count=0}이기 때문에 유효한 값이 하나도 없어도 덧셈의 항등원인 0을 반환할 수 있다.

반면 \zsh{min_count=1}을 지정하면 최소 한 개의 유효한 값이 있어야 집계 결과를 인정한다.

이 경우 유효한 값이 하나도 없으므로 \zsh{<NA>}가 반환된다.

따라서 집계 결과가 0이라는 사실만으로 실제 관측값의 합이 0이었다고 결론 내리면 안 된다.

이것은 학습 로그나 평가 지표를 집계할 때 특히 중요한 경계 조건이다.

\subsubsection{ML/Transforemr 응용 - 누락된 레이블과 잘못된 전처리 결과 구분}

이제 이번 학습의 내용을 실제 데이터 분석 파이프라인에 연결해 본다.

문서 분류 모델을 학습하기 위한 데이터셋이 있다고 가정한다.

각 문서에는 텍스트, 정답 레이블, 토큰 길이, 데이터 분할 정보가 있다.

그런데 일부 데이터에 문제가 발생했다.

\begin{dhksmry}
    - 문서 d2 : 원본 텍스트 누락

    - 문서 d3 : 정답 레이블 누락, 토큰 길이 변환 실패

    - 문서 d4 : 정답 레이블 형식 오류
\end{dhksmry}

이 상황을 Pandas 객체로 표현해 보자.

\begin{dhknote}
\begin{dhksmry}[centered]
독립 실행 코드 7 - 데이터 품질 상태를 보존하는 전처리
\end{dhksmry}
\end{dhknote}

\begin{dhkcode}
import pandas as pd

docs = pd.DataFrame(
    {
        "text": [
            "good",
            None,
            "bad",
            "ok",
        ],
        "raw_label": [
            "1",
            "0",
            None,
            "unknown",
        ],
        "raw_token_length": [
            "5",
            "3",
            "bad",
            "8",
        ],
        "split": [
            "train",
            "train",
            "valid",
            "valid",
        ],
    },
    index=pd.Index(
        ["d1", "d2", "d3", "d4"],
        name="doc_id",
    ),
)

docs["label"] = pd.to_numeric(
    docs["raw_label"],
    errors="coerce",
).astype("Int64")

docs["token_length"] = pd.to_numeric(
    docs["raw_token_length"],
    errors="coerce",
).astype("Int64")

docs["invalid_label"] = (
    docs["raw_label"].notna()
    & docs["label"].isna()
)

docs["invalid_token_length"] = (
    docs["raw_token_length"].notna()
    & docs["token_length"].isna()
)

docs["usable"] = (
    docs["text"].notna()
    & docs["label"].notna()
    & docs["token_length"].notna()
)

print(
    docs[
        [
            "label",
            "token_length",
            "invalid_label",
            "invalid_token_length",
            "usable",
        ]
    ]
)

print()
print(docs.index[docs["usable"]].tolist())
\end{dhkcode}

예상 출력:

\begin{dhkzsh}
        label  token_length  invalid_label  invalid_token_length  usable
doc_id
d1          1             5          False                 False    True
d2          0             3          False                 False   False
d3       <NA>          <NA>          False                  True   False
d4       <NA>             8           True                 False   False

['d1']
\end{dhkzsh}

\begin{dhknote}
\begin{dhksmry}[centered]
객체와 실행 과정 분석
\end{dhksmry}
\end{dhknote}

\begin{dhknote}
\begin{dhksmry}[centered]
1단계 - 원본 DataFrame 생성
\end{dhksmry}
\end{dhknote}

\zsh{docs}라는 이름은 DataFrame 객체에 바인딩된다.

이 DataFrame은 원본 문자열과 결측 상태를 그대로 보관하고 있다.

\begin{dhknote}
\begin{dhksmry}[centered]
2단계 - 정답 레이블 반환
\end{dhksmry}
\end{dhknote}

\begin{dhkcode}
docs["label"] = ...
\end{dhkcode}

\zsh{pd.to_numeric()}이 숫자 변환을 수행하고, 

\zsh{astype("Int64")}가 Nullable 정수형으로 변환한다.

이 결과는 DataFrame의 새로운 열로 할당된다.

이때 원본 \zsh{raw-label} 열은 삭제하거나 덮어쓰지 않았다.

\begin{dhknote}
\begin{dhksmry}[centered]
3단계 - 파싱 실패 추적
\end{dhksmry}
\end{dhknote}

\begin{dhkcode}
docs["invalid_label"] = (
    docs["raw_label"].notna()
    & docs["label"].isna()
)
\end{dhkcode}

원본에는 값이 있었지만 숫자로 변환되지 못한 경우를 별도로 표시한다.

따라서 \zsh{d3}의 누락된 레이블과 \zsh{d4}의 잘못된 레이블을 구분할 수 있다.

\begin{dhknote}
\begin{dhksmry}[centered]
4단계 - 사용 가능 여부 판단
\end{dhksmry}
\end{dhknote}

\begin{dhkcode}
docs["usable"] = (
    docs["text"].notna()
    & docs["label"].notna()
    & docs["token_length"].notna()
)
\end{dhkcode}

이 표현식은 세 가지 조건을 모두 만족하는 문서를 선택하기 위한 불리언 Series를 생성한다.

현재 예제에서는 \zsh{d1}만 사용 가능한 것으로 판정된다.

다만 이 코드는 모든 프로젝트에 적용되는 보편적인 데이터 품질 정책은 아니다.

실제 Transformer 학습에서는 토큰 길이를 다시 계산할 수 있을 수도 있고, 텍스트가 누락된 문서는 재수집 대상이 될 수도 있다.

중요한 것은 정제 규칙을 적용하기 전에 어떤 문제가 발생했는지 보존했다는 점이다.

\begin{dhknote}
\begin{dhksmry}[centered]
PyTorch로 전달할 때 주의할 점
\end{dhksmry}
\end{dhknote}

Pandas의 Nullable dtype은 데이터 분석에 편리하지만, PyTorch Tensor로 직접 전달할 수 있는 자료형과 항상 일치하지는 않는다.

예를 들어 \zsh{Int64} Series에 \zsh{<NA>}가 포함되어 있다면, Tensor를 생성하기 전에 결측을 처리해야 한다.

또한 \zsh{.to_numpy()}로 변환하는 과정에서 dtype과 결측 표현이 바뀔 수 있으므로 변환 결과를 확인해야 한다.

\begin{dhkcode}
import pandas as pd

s = pd.Series(
    [1, pd.NA],
    dtype="Int64",
)

print(s.to_numpy(dtype="int64", na_value=-1))
\end{dhkcode}

실행 결과:

\begin{dhkzsh}
[ 1 -1]
\end{dhkzsh}

이 코드는 결측을 \zsh{-1}로 치환했다.

그러나 \zsh{-1}이 실제로 유효한 레이블이라면 결측과 유효한 값을 구분할 수 없게 된다.

따라서 모델 입력으로 변환할 때는 다음 세 조건을 함께 확인해야 한다.

\begin{dhksmry}
1. 결측값이 어떤 정책으로 처리되었는가?

2. 변환된 배열의 dtype이 모델이 기대하는 자료형인가?

3. 변환 후에도 문서 ID와 행 순서가 정답 레이블에 정확하게 대응하는가?
\end{dhksmry}

특히 세 번째 조건은 \zsh{D.5}와 \zsh{D.21}에서 본격적으로 다룰 것이다.

\subsubsection{자주 발생하는 오해와 경계 조건}

이번 학습의 핵심 오해를 비교해 보자.

\begin{dhktable}
|오해	|정확한 이해|
|:---|:---|
|None과 np.nan은 같은 객체다|	서로 다른 타입과 비교 규칙을 갖는 객체다
|np.nan == np.nan으로 결측을 검사할 수 있다|	NaN의 자기 비교는 거짓이므로 isna()를 사용해야 한다
|pd.NA는 항상 False처럼 동작한다|	진리값이 모호하며, 일부 논리 연산에서만 결과를 확정할 수 있다
|int64와 Int64는 표기만 다르다|	NumPy 정수 dtype과 Pandas Nullable 정수 dtype이다
|object 열은 문자열 열이다|	Python 객체를 담는 범용 자료형이다
|errors="coerce"는 잘못된 데이터를 수정한다|	변환 실패를 결측으로 표현할 뿐, 원인을 자동으로 보존하지 않는다
|sum()이 0이면 합계가 0으로 관측되었다|	유효한 관측값이 없어도 기본 설정에서 0이 반환될 수 있다
|Nullable dtype은 모든 내부 구현이 같다|	자료형과 백엔드에 따라 저장 방식이 달라진다
\end{dhktable}

여기에 버전 의존성도 고려해야 한다.

특히 Pandas 2.x와 3.x는 문자열 자료형의 기본 추론, Copy-on-Write 동작 등에서 차이가 있다.

따라서 학습 코드에서는 \zsh{dtype}을 명시하고, 결과를 재현할 때 Pandas 버전을 함께 기록하는 것이 좋다.

이번 학습의 코드는 로컬 환경의 Pandas 2.2.3, NumPy 2.3.5에서 실행해 주요 출력과 오류 동작을 확인하였다.

\subsubsection{핵심 정리 및 다음 학습}

이번 \zsh{D.3}에서 가장 중요한 결론은 세 가지다.

첫째, Python 객체 타입과 데이터 저장 자료형은 다른 개념이다.

\zsh{type(s)}는 Series 객체의 타입을 나타내고, \zsh{s.dtype}은 Series가 관리하는 값들의 자료형을 나타낸다.

둘째, 결측값은 하나의 공통 객체로 통일되어 있지 않다.

\zsh{None}, \zsh{np.nan}, \zsh{pd.NA}, \zsh{pd.NaT}는 서로 다른 객체와 연산 규칙을 갖지만 Pandas의 결측 검사 인터페이스를 통해 공통적으로 다룰 수 있다.

셋째, 결측 처리에는 데이터의 의미에 관한 판단이 필요하다.

\zsh{fillna(0)}, \zsh{errors="coerce"}, \zsh{dropna()}는 단순한 코드 조작이 아니라 데이터의 의미와 분석 대상 자체를 바꿀 수 있는 결정이다.

\zsh{D.3}에서는 Series가 어떤 dtype으로 값을 관리하는지 배웠다.

이제 다음 질문이 자연스럽게 발생할 것이다.

\begin{dhkcode}
s = pd.Series(
    [10, 20, 30],
    index=[2, 0, 1],
)
\end{dhkcode}

이때 \zsh{s[0]}은 첫 번째 원소인 10을 반환할까, 아니면 레이블 0에 대응하는 20을 반환할까?

\zsh{D.4}에서는 \zsh{[]}, \zsh{.loc}, \zsh{.iloc}, \zsh{.at}, \zsh{.iat}의 동작을 객체의 속성 탐색과 인덱싱 프로토콜에 연결해 분석할 것이다.

특히 정수 레이블과 정수 위치가 충돌하는 상황, 반환 타입, 슬라이스 경계, 중복 레이블의 영향을 다룰 것이다.

공식 참고 자료:

\begin{dhksmry}
https://pandas.pydata.org/docs/user_guide/missing_data.html
https://pandas.pydata.org/docs/user_guide/integer_na.html
https://pandas.pydata.org/docs/user_guide/boolean.html
https://pandas.pydata.org/docs/user_guide/categorical.html
\end{dhksmry}
